% ============================================================================
\section{Introduction}
\label{sec:intro}

\subsection{Motivation}

About 50 million people live with epilepsy. Roughly one in three has drug-resistant seizures, and people with frequent tonic-clonic seizures carry a markedly raised risk of sudden unexpected death in epilepsy (SUDEP) that is reduced by nocturnal supervision and seizure alarms (Section~\ref{sec:background}). Patients also under-report their own seizures, so objective seizure counts are clinically valuable even when no alarm is raised. Existing cleared wearables (wrist accelerometry and electrodermal activity, upper-arm PPG, surface EMG) work well for \emph{convulsive} seizures but largely miss focal seizures, which are the majority.

In 2025--2026, Deepinder Goyal (founder of Zomato) presented \emph{Temple}, a small temple-worn optical wearable. This project began by asking what exactly Temple measures and whether a device of that size and at that location could detect seizures. The answer led, step by step, to a different device concept, a ``?''-shaped ear-hook worn on both ears that combines ear-EEG, in-ear PPG and motion sensing, and to a sequence of simulations and data studies to test it.

\subsection{Research questions}

The project was organised around a chain of questions, each answered by a document, a piece of code or an experiment (Table~\ref{tab:journey}).

\begin{table}[H]
\centering\small
\caption{Research journey: questions, methods and outputs.}
\label{tab:journey}
\begin{tabularx}{\linewidth}{@{}c L L@{}}
\toprule
\# & Question & Output (section) \\
\midrule
1 & What exactly is Temple, technically? & Technical dossier (\S\ref{sec:temple}) \\
2 & What is epilepsy, and what happens in the brain during a seizure? & Background review (\S\ref{sec:background}) \\
3 & Can seizures be simulated realistically, including ion concentrations? & Three simulators; ionic model validated against Cressman 2009 (\S\ref{sec:sim}) \\
4 & Which hidden variables can a wearable infer, and with what formulas? & Derivations and verification audit (\S\ref{sec:inference}) \\
5 & Do the formulas work? Does our code match an established simulator? & Twin experiments E1--E7; TVB cross-validation (\S\ref{sec:twin}) \\
6 & How well do temple or ear electrodes detect seizures on real EEG? & CHB-MIT patient-specific study (\S\ref{sec:chbmit}) \\
7 & Can one model work for new users? How much does personal fine-tuning help? & CHB-MIT general model, 23 subjects (\S\ref{sec:chbmit}) \\
8 & What should the device look like? & ``?'' ear-hook concept (\S\ref{sec:device}) \\
9 & Can the whole system be made visible and intuitive? & Offline interactive Live Lab (\S\ref{sec:livelab}) \\
10 & Which physiological signals change during seizures, and which are measurable at the ear? & Body-signal review (\S\ref{sec:bodysignals}) \\
11 & Do heart, muscle and motion signals cut false alarms on real wearable data? & SeizeIT2 multimodal study, 125 patients (\S\ref{sec:seizeit2}) \\
12 & How far can personalisation, context and alarm logic push performance, and what limits it? & 22 SeizeIT2 experiments (\S\ref{sec:seizeit2}) \\
\bottomrule
\end{tabularx}
\end{table}

\subsection{Contributions}

\begin{enumerate}
\item A structured technical analysis of Temple (confirmed, reported and inferred facts separated), concluding that it cannot detect seizures as shipped.
\item A validated biophysical seizure simulator (\code{ionic.py}) with \Kp, \Nap, \Clm, KCC2, glia, an electrogenic pump, a \Nap-activated \Kp{} current, adenosine and Balloon--Windkessel haemodynamics, documented with a 14-step calibration log.
\item Derivations and a verification audit of inference formulas for hidden seizure variables, with three errors in earlier drafts caught and corrected, and twin experiments that confirm when each formula works and fails.
\item A numerical cross-validation of our Epileptor against The Virtual Brain and a VEP-style whole-brain inference, with an explicit discussion of the ``inverse crime''.
\item Detection studies on real scalp EEG (CHB-MIT) comparing full cap, temple and ear-level montages, per-person normalisation, and personal fine-tuning.
\item An offline, single-screen visualisation (``Live Lab'') that couples the simulator to a virtual ear-EEG, a body-signal model and EEG-only vs multimodal detectors.
\item A large-scale multimodal study on SeizeIT2 (125 patients): streaming 126\,GB through memory, feature engineering for ear-EEG/ECG/EMG/motion, and 22 experiments on fusion, context, a novel movement-adjusted heart-rate feature, personalisation, two-stage verification, simulated in-ear PPG, seizure-type analysis and a sensitivity-ceiling diagnosis.
\item A transparent record of every failure, bug and correction.
\end{enumerate}

\subsection{Tools and environment}

All code is Python~3.13 managed with \code{uv}, using numpy, scipy, numba, matplotlib, mne, scikit-learn, fooof and tqdm; The Virtual Brain (\code{tvb-library}, \code{tvb-data}) and CPU PyTorch are optional extras. Work ran on a Windows~11 laptop (12 logical cores, 15.3\,GB RAM) with limited free disk, which shaped several engineering choices (streaming, feature-only storage, detached background jobs). The repository is \code{rhydham-mittal27/epilepsy-temple} on GitHub (private).
